\documentclass[11pt]{article}
\usepackage[a4paper,margin=1in]{geometry}
\usepackage{fontspec}
\usepackage[boldfont=false]{xeCJK}
\setCJKmainfont{FandolSong-Regular.otf}
\usepackage{amsmath}
\usepackage{amssymb}
\usepackage{array}
\usepackage{booktabs}
\usepackage{graphicx}
\usepackage{adjustbox}
\usepackage{enumitem}
\setlist[itemize]{topsep=2pt,partopsep=0pt,parsep=0pt,itemsep=2pt,leftmargin=1.4em}
\setlist[enumerate]{topsep=2pt,partopsep=0pt,parsep=0pt,itemsep=2pt,leftmargin=1.6em}
\usepackage{parskip}
\usepackage{fvextra}
\fvset{breaklines=true,breakanywhere=true,fontsize=\small}
\setlength{\parindent}{0pt}

\begin{document}

\begin{center}
  {\LARGE\bfseries Enzymes}\\[0.35em]
  {\large A Level Biology}
\end{center}
\vspace{0.6em}

\section*{What enzymes are}

\textbf{Enzymes} 酶 are \textbf{globular proteins} 球状蛋白质 — a type of \textbf{protein} 蛋白质 with a rounded, soluble shape. They are biological catalysts: they \textbf{catalyse} 催化 (speed up) the chemical reactions in living things, and they are not used up, so each enzyme works again and again.

Enzymes work in two places:

\begin{itemize}
\item \textbf{intracellular} 细胞内 enzymes work inside the \textbf{cell} 细胞 that made them. An example is \textbf{catalase} 过氧化氢酶, which breaks down harmful hydrogen peroxide.

\item \textbf{extracellular} 细胞外 enzymes are \textbf{secreted} 分泌 (sent out) to work outside the cell. An example is \textbf{amylase} 淀粉酶, which is released into the gut to \textbf{digest} 消化 \textbf{starch} 淀粉.

\end{itemize}

\section*{How enzymes work}

Each enzyme has a special pocket called the \textbf{active site} 活性位点. The molecule it acts on is its \textbf{substrate} 底物. The substrate fits into the active site to form an \textbf{enzyme–substrate complex} 酶底物复合物. The reaction then happens, and the \textbf{products} 产物 leave, freeing the active site for the next substrate.

\subsection*{Specificity}

An enzyme is \textbf{specific}: it usually works on only one substrate. This is because the shape of the active site is \textbf{complementary} 互补 to (fits) the shape of that substrate and no other. We call this \textbf{specificity} 专一性.

Two ideas explain how the substrate fits:

\begin{itemize}
\item the \textbf{lock-and-key hypothesis} 锁钥学说 — the active site is a fixed shape, and only a substrate with the matching shape fits, like a key in a lock.

\item the \textbf{induced-fit hypothesis} 诱导契合学说 — the active site is not quite the right shape at first. When the substrate binds, the active site changes shape a little to wrap around it tightly. This idea fits the evidence better.

\end{itemize}

\subsection*{Lowering activation energy}

Every reaction needs a small "push" of energy to start, called the \textbf{activation energy} 活化能. An enzyme lowers the activation energy. This lets the reaction go quickly at the cell's normal \textbf{temperature} 温度, instead of needing high heat.

\section*{Measuring the rate of a reaction}

You can follow an enzyme reaction in two ways:

\begin{itemize}
\item measure how fast \textbf{product} is made. With catalase, oxygen gas is a product, so you collect the gas and measure its volume over time.

\item measure how fast substrate disappears. With amylase, you remove samples and use the iodine test; the blue-black colour fades as the starch is used up.

\end{itemize}

A \textbf{colorimeter} 比色计 makes this exact. It shines light through the tube and measures how much light is absorbed, so a colour change becomes a number you can plot.

The \textbf{rate of reaction} 反应速率 is steepest at the start (most substrate present), so the \textbf{initial rate} (the slope at time zero) is the fairest value to compare.

\section*{Factors that affect enzyme activity}

\subsection*{Temperature}

As temperature rises, molecules gain more \textbf{kinetic energy} 动能 and \textbf{collide} 碰撞 more often, so the rate rises. But above the \textbf{optimum temperature} 最适温度 the enzyme begins to \textbf{denature} 变性: the heat breaks the bonds holding its shape, so the active site changes and no longer fits the substrate. The rate then falls quickly.

\subsection*{pH}

Each enzyme has an optimum pH. If the pH moves too far from it, the enzyme denatures and the rate drops. To study pH fairly, you keep it steady with a \textbf{buffer solution} 缓冲液.

\subsection*{Enzyme concentration}

With plenty of substrate, more enzyme means more active sites, so the rate goes up in proportion to enzyme concentration.

\subsection*{Substrate concentration}

At first, adding more substrate speeds the reaction. But once every active site is busy, adding more makes no difference — the rate levels off at a maximum.

\subsection*{Inhibitor concentration}

An \textbf{inhibitor} 抑制剂 is a molecule that slows an enzyme. The more inhibitor present, the lower the rate.

\section*{V\_max and the Michaelis–Menten constant}

The levelling-off rate, when all active sites are full, is the maximum rate, written $V_{\text{max}}$.

The \textbf{Michaelis–Menten constant} 米氏常数 ($K_{\text{m}}$) is the substrate concentration that gives \textbf{half} of $V_{\text{max}}$. It tells you about the enzyme's \textbf{affinity} 亲和力 (pulling power) for its substrate:

\begin{itemize}
\item a \textbf{low} $K_{\text{m}}$ means the enzyme reaches half-speed at a low substrate concentration, so it has a \textbf{high} affinity.

\item a \textbf{high} $K_{\text{m}}$ means a \textbf{low} affinity.

\end{itemize}

So $K_{\text{m}}$ lets you compare how strongly different enzymes hold their substrates.

\section*{Reversible inhibitors}

Some inhibitors are \textbf{reversible} 可逆: they can leave the enzyme again. There are two types.

\begin{center}
\adjustbox{max width=\linewidth}{%
\begin{tabular}{llll}
\toprule
\textbf{Type} & \textbf{Where it binds} & \textbf{Effect of adding more substrate} & \textbf{Effect on $V_{\text{max}}$ and $K_{\text{m}}$} \\
\midrule
\textbf{competitive inhibitor} 竞争性抑制剂 & in the active site (it has a similar shape to the substrate) & more substrate out-competes it, so its effect is reduced & $V_{\text{max}}$ unchanged; $K_{\text{m}}$ rises \\
\textbf{non-competitive inhibitor} 非竞争性抑制剂 & at another site, changing the active site's shape & adding more substrate does not help & $V_{\text{max}}$ falls; $K_{\text{m}}$ unchanged \\
\bottomrule
\end{tabular}}
\end{center}

\section*{Immobilised enzymes}

An \textbf{immobilised enzyme} 固定化酶 is fixed in place — for example, trapped inside small beads of \textbf{alginate} 海藻酸盐 — instead of floating free in solution. The substrate solution flows past the beads.

A free enzyme usually works a little faster, because the substrate can reach it easily. But immobilised enzymes have big practical advantages:

\begin{itemize}
\item the enzyme is not washed away, so it can be used again and again.

\item the product is pure — it is not mixed with enzyme.

\item the enzyme is more stable, so it survives changes in temperature and pH better.

\item the process can run continuously, with substrate flowing in and product flowing out.

\end{itemize}


\end{document}
